\section{Chapter 3: Functions, definitions, and structures}\label{sec:15-appendix-solutions:03-chapter-3:1}

\textbf{1. Why \texttt{⟨0, 0⟩} is unambiguous}

\texttt{def origin : Point := ⟨0, 0⟩} type-checks unambiguously because the \emph{expected type} \texttt{Point} is already fixed by the \texttt{def}'s signature before the right-hand side is elaborated. Lean reads off \texttt{Point}'s constructor and field order from that expected type alone, so \texttt{⟨0, 0⟩} needs no further information to know it means \texttt{Point.mk 0 0} rather than, say, \texttt{Pair.mk 0 0}.

This property fails whenever no expected type is available, or several structures with the same field arity are simultaneously in scope with nothing to disambiguate them: \texttt{\#check (⟨0, 0⟩ : \_\allowbreak{})} with the placeholder left fully open. Lean cannot choose a constructor from field values alone, and reports an ``expected type'' error rather than a fabricated guess. The named form \texttt{\{ x := ..., y := ... \}} sidesteps this, since the field names themselves pin down the constructor regardless of context.

\textbf{2. Writing out \texttt{.toPoint} by hand}

Without \texttt{extends}, the projection \texttt{extends Point} generates would have to be written explicitly, as \texttt{def Point3D.toPoint (p : Point3D) : Point := \{ x := p.x, y := p.y \}}.

This is precisely a \textbf{forgetful functor} (\hyperref[sec:02-terminology-and-coc:01-terminology:1]{Chapter 2, Section 1}, and the Mathematical reading box of \hyperref[sec:03-functions-and-structures:03-extending-structures:1]{Section 3}): a map that keeps some of the data of a structure (here, \texttt{x} and \texttt{y}) and discards the rest (\texttt{z}). \texttt{extends} generates exactly this projection automatically, under the name \texttt{.toPoint}, rather than requiring it to be written by hand for every extension.

\textbf{3. Proof fields change nothing about type-checking}

Claim: a \texttt{structure} field whose declared type is a proposition \texttt{P} is checked by exactly the same mechanism as a field whose declared type is \texttt{Nat} or any other \texttt{Type}.

\emph{Proof.} Field checking in a \texttt{structure} constructor is uniform across all fields: given \texttt{\{ f₁ := e₁, ..., fₙ := eₙ \}}, Lean elaborates each \texttt{eᵢ} against the declared type of \texttt{fᵢ} and accepts the term exactly when that type-check succeeds, regardless of which universe (\texttt{Type} or \texttt{Prop}) the declared type inhabits. Supplying a proof field is the ordinary case of this rule with \texttt{fᵢ : P} for some \texttt{P : Prop}; the term \texttt{eᵢ} supplied must be a proof of \texttt{P}, checked once, at construction time, the same way \texttt{eᵢ} supplied for a \texttt{Nat}-typed field must genuinely have type \texttt{Nat}. \(\blacksquare\)

This is exactly what makes \texttt{Group} (Chapter 7) impossible to build carelessly. Its axiom fields (\texttt{assoc}, \texttt{id\_\allowbreak{}left}, \ldots) are \texttt{Prop}-valued fields like any other, so a \texttt{Group} value cannot be assembled without supplying an actual proof of each axiom, checked by the same constructor mechanism that checks \texttt{op}, \texttt{id}, and \texttt{inv}.

\textbf{4. \texttt{Rectangle} and \texttt{area}}

\begin{lstlisting}
structure Rectangle where
  width : Nat
  height : Nat

def area (r : Rectangle) : Nat := r.width * r.height

#eval area (*$\langle$*)3, 4(*$\rangle$*)   -- 12
\end{lstlisting}

\texttt{⟨3, 4⟩} is the anonymous constructor, filling \texttt{width} and \texttt{height} in field-declaration order. \texttt{area} reads both fields back out via \texttt{.width} and \texttt{.height} projection and multiplies them. Nothing here differs from \texttt{Point}/\texttt{shift} in Section 1, only the field names and the arithmetic change.

\textbf{5. \texttt{Box α} and \texttt{unwrap}}

\begin{lstlisting}
structure Box ((*$\alpha$*) : Type) where
  value : (*$\alpha$*)

def unwrap {(*$\alpha$*) : Type} (b : Box (*$\alpha$*)) : (*$\alpha$*) := b.value

def natBox : Box Nat := (*$\langle$*)7(*$\rangle$*)
def strBox : Box String := (*$\langle$*)"seven"(*$\rangle$*)

#eval unwrap natBox   -- 7
#eval unwrap strBox   -- "seven"
\end{lstlisting}

\texttt{Box} is \texttt{Pair} with one slot instead of two, the same ``parameterize by the type itself'' idea from Section 2, at its simplest possible size. \texttt{Box Nat} and \texttt{Box String} are two different types, generated from the one \texttt{structure Box (α : Type)} declaration, the way \texttt{Pair Nat String} and \texttt{Pair String Nat} are two different instantiations of \texttt{Pair}. \texttt{unwrap} is implicit in \texttt{α}, exactly like \texttt{identity} in Chapter 1, since the type of \texttt{b.value} already determines it.

\textbf{6. \texttt{ColoredRectangle} and the forgetful projection}

\begin{lstlisting}
structure ColoredRectangle extends Rectangle where
  color : String

def redSquare : ColoredRectangle :=
  { width := 5, height := 5, color := "red" }

#check redSquare.toRectangle      -- redSquare.toRectangle : Rectangle
#eval area redSquare.toRectangle  -- 25
\end{lstlisting}

\texttt{redSquare.toRectangle} has type \texttt{Rectangle}, generated automatically by \texttt{extends}, discarding only the \texttt{color} field and keeping \texttt{width}/\texttt{height} untouched. \texttt{area} was written once, against \texttt{Rectangle}, with no knowledge that \texttt{ColoredRectangle} would ever exist. Calling it on \texttt{redSquare.toRectangle} reuses that same function unchanged, computing \texttt{25} from the inherited \texttt{width}/\texttt{height} fields alone. This is exactly the forgetful-functor pattern from the Mathematical reading box of Section 3, not a coincidence of naming. \texttt{.toRectangle} is the map \(U : \mathrm{ColoredRectangle} \to \mathrm{Rectangle}\) that strips away the extra structure (the \texttt{color} field) and keeps everything else, and \texttt{area} factoring through it is exactly what it means for \texttt{area} to be a function defined on the \emph{underlying} \texttt{Rectangle}, indifferent to whatever extra data the specific extension of a caller happens to add on top.
